\chapter*{Chapitre 9 — Acoustique du bâtiment}
\addcontentsline{toc}{chapter}{Chapitre 9 — Acoustique du bâtiment}

% ============================================================
\section*{Énoncés}
% ============================================================

\exercice{Calcul de niveaux acoustiques}

Un extracteur de cuisine collective est caractérisé par un niveau de puissance sonore $L_w = 88$~dB. Il est installé au centre d'un local de dimensions $8 \times 6 \times 3$~m (coefficient d'absorption moyen $\bar{\alpha} = 0{,}15$, source au plafond $Q = 2$).

\begin{enumerate}
  \item Calculer la constante du local $R$.
  \item Calculer le niveau de pression sonore $L_p$ à 1,5~m de l'extracteur.
  \item Quel serait $L_p$ si l'on passait à $\bar{\alpha} = 0{,}40$ grâce à un traitement acoustique des parois ?
  \item Deux extracteurs identiques sont mis en service simultanément. Quel est le niveau $L_p$ résultant à 1,5~m (avec $\bar{\alpha} = 0{,}15$) ?
\end{enumerate}

\textbf{Données :}
\[
  L_p = L_w + 10 \cdot \log_{10}\!\left(\frac{Q}{4\pi r^2} + \frac{4}{R}\right), \qquad R = \frac{S \cdot \bar{\alpha}}{1 - \bar{\alpha}}
\]

\exercice{Isolement acoustique — choix de paroi}

Un local de compresseurs (niveau de pression mesuré $L_1 = 92$~dB(A)) jouxte un bureau ouvert (exigence $L_2 \leq 45$~dB(A)). La salle de réception a une absorption $A = 25$~sabins et la paroi séparative mesure $S = 20$~\si{\metre\squared}.

\begin{enumerate}
  \item Calculer l'affaiblissement acoustique $R$ minimum de la paroi séparative.
  \item En déduire la masse surfacique minimale $m_s$ à 500~Hz (loi de masse).
  \item Proposer une solution constructive et vérifier sa conformité à partir du tableau de $R_w$ ci-dessous.
  \item Citer deux mesures complémentaires pour limiter la transmission solidienne.
\end{enumerate}

\textbf{Tableau de référence $R_w$ :}
\begin{center}
\begin{tabular}{lc}
\toprule
Type de paroi & $R_w$ (dB) \\
\midrule
Cloison légère plaque de plâtre 72 mm & 38 \\
Mur béton 15 cm & 48 \\
Mur béton 20 cm & 52 \\
Porte acoustique renforcée & 42--47 \\
\bottomrule
\end{tabular}
\end{center}

\textbf{Formules utiles :}
\[
  R = L_1 - L_2 + 10 \cdot \log_{10}\!\left(\frac{S}{A}\right), \qquad R \approx 20 \cdot \log_{10}(m_s \cdot f) - 48
\]

\exercice{Dimensionnement acoustique d'une installation VMC double flux}

Une CTA double flux dessert 8 bureaux individuels (exigence NC 30 $\approx$ 38~dB(A)). Les données constructeur de la CTA sont : $L_w = 78$~dB (global) ; le réseau aéraulique assure une atténuation de 15~dB ; les bouches de soufflage ont un niveau NC 28 pour le débit nominal.

\begin{enumerate}
  \item Calculer l'atténuation totale nécessaire entre la CTA et les locaux.
  \item Déterminer si un silencieux est nécessaire et quelle atténuation il doit apporter.
  \item La vitesse dans la gaine principale est de 5~m/s pour un débit de 800~m³/h. Calculer la section de gaine rectangulaire ($a \times b$ avec $a = 400$~mm). Cette vitesse est-elle conforme aux recommandations ?
  \item Citer deux réglages opérationnels permettant de réduire le bruit sans modifier l'installation.
\end{enumerate}

\textbf{Rappel :} $Q = v \times S_{\text{gaine}}$, avec $Q$ en m³/s et $S$ en m².

\newpage
% ============================================================
\section*{Corrigés}
% ============================================================

\subsection*{Corrigé — Exercice 1}

\textit{Question 1 — Constante du local :}

Surface totale des parois : $S = 2 \times (8 \times 6 + 8 \times 3 + 6 \times 3) = 2 \times (48 + 24 + 18) = 180$~m²
\[
  R = \frac{S \cdot \bar{\alpha}}{1 - \bar{\alpha}} = \frac{180 \times 0{,}15}{1 - 0{,}15} = \frac{27}{0{,}85} = \boxed{31{,}8~\text{m}^2}
\]

\textit{Question 2 — Niveau de pression à 1,5~m ($\bar{\alpha} = 0{,}15$) :}
\[
  \frac{Q}{4\pi r^2} = \frac{2}{4\pi \times 1{,}5^2} = \frac{2}{28{,}27} = 0{,}0707
\]
\[
  \frac{4}{R} = \frac{4}{31{,}8} = 0{,}1258
\]
\[
  L_p = 88 + 10 \cdot \log_{10}(0{,}0707 + 0{,}1258) = 88 + 10 \cdot \log_{10}(0{,}1965) = 88 - 7{,}1 = \boxed{80{,}9~\text{dB}}
\]

\textit{Question 3 — Avec $\bar{\alpha} = 0{,}40$ :}
\[
  R' = \frac{180 \times 0{,}40}{1 - 0{,}40} = \frac{72}{0{,}60} = 120~\text{m}^2
\]
\[
  \frac{4}{R'} = \frac{4}{120} = 0{,}0333
\]
\[
  L_p' = 88 + 10 \cdot \log_{10}(0{,}0707 + 0{,}0333) = 88 + 10 \cdot \log_{10}(0{,}1040) = 88 - 9{,}8 = \boxed{78{,}2~\text{dB}}
\]

\begin{attention}
Le gain acoustique par traitement absorbant est de $80{,}9 - 78{,}2 = 2{,}7$~dB seulement à 1,5~m. Le champ direct domine à courte distance ; le traitement absorbant est surtout efficace en champ réverbéré (grande distance).
\end{attention}

\textit{Question 4 — Deux extracteurs identiques ($L_w = 88$~dB chacun) :}

Le niveau de puissance sonore de deux sources identiques :
\[
  L_{w,total} = 88 + 10 \cdot \log_{10}(2) = 88 + 3 = 91~\text{dB}
\]
\[
  L_p = 91 + 10 \cdot \log_{10}(0{,}0707 + 0{,}1258) = 91 - 7{,}1 = \boxed{83{,}9~\text{dB}}
\]

\subsection*{Corrigé — Exercice 2}

\textit{Question 1 — Affaiblissement acoustique $R$ minimum :}
\[
  R_{\min} = L_1 - L_2 + 10 \cdot \log_{10}\!\left(\frac{S}{A}\right) = 92 - 45 + 10 \cdot \log_{10}\!\left(\frac{20}{25}\right)
\]
\[
  = 47 + 10 \cdot \log_{10}(0{,}80) = 47 - 1{,}0 = \boxed{46~\text{dB}}
\]

\textit{Question 2 — Masse surfacique minimale à 500~Hz :}

À partir de la loi de masse $R \approx 20 \cdot \log_{10}(m_s \cdot f) - 48$ :
\[
  46 = 20 \cdot \log_{10}(m_s \times 500) - 48
\]
\[
  20 \cdot \log_{10}(500 \cdot m_s) = 94
\]
\[
  \log_{10}(500 \cdot m_s) = 4{,}7 \quad \Rightarrow \quad 500 \cdot m_s = 10^{4{,}7} = 50\,119
\]
\[
  m_s = \frac{50\,119}{500} = \boxed{100~\text{kg/m}^2}
\]

\textit{Question 3 — Solution constructive :}

Un mur béton de 15~cm ($R_w = 48$~dB) répond à l'exigence de $R_{\min} = 46$~dB. Un mur béton de 20~cm ($R_w = 52$~dB) offre une marge de sécurité supplémentaire de 6~dB.

\begin{methode}
Vérification : $R_w = 48$~dB $>$ $R_{\min} = 46$~dB \checkmark \quad La paroi béton 15~cm est conforme.
\end{methode}

\textit{Question 4 — Mesures complémentaires pour la transmission solidienne :}
\begin{itemize}
  \item Montage du compresseur sur plots anti-vibratiles ou ressorts isolants pour découpler les vibrations du bâti.
  \item Raccordement des canalisations par manchons souples (EPDM) pour éviter la propagation des vibrations vers les parois.
\end{itemize}

\subsection*{Corrigé — Exercice 3}

\textit{Question 1 — Atténuation totale nécessaire :}

L'exigence NC 30 correspond à environ 38~dB(A) en réception. L'atténuation totale nécessaire entre la CTA et les locaux est :
\[
  \text{Att}_{totale} = L_{w,CTA} - L_{p,exigence} = 78 - 38 = 40~\text{dB}
\]

\textit{Question 2 — Silencieux nécessaire :}

L'atténuation naturelle du réseau est de 15~dB. La bouche de soufflage émet NC 28 $\approx$ 36~dB(A), ce qui est inférieur à 38~dB(A) — elle est conforme. Mais il faut vérifier le niveau en sortie de réseau avant la bouche :
\[
  L_{p,\text{avant bouche}} = 78 - 15 = 63~\text{dB}
\]

Atténuation manquante :
\[
  \text{Att}_{\text{silencieux}} = 63 - 38 = \boxed{25~\text{dB}}
\]

Un silencieux à baffles de 1~m (atténuation $\approx 25$~dB(A) à 500~Hz) est nécessaire. Vérifier la perte de charge ajoutée ($\approx$ 35~Pa) par rapport à la pression disponible du ventilateur.

\textit{Question 3 — Section de gaine et vitesse :}

Débit : $Q = \SI{800}{\m^3 \per \hour} = \SI{800}{} / 3600 = 0{,}222~\text{m}^3/\text{s}$

Section nécessaire :
\[
  S_{\text{gaine}} = \frac{Q}{v} = \frac{0{,}222}{5} = 0{,}0444~\text{m}^2
\]

Avec $a = 400$~mm $= 0{,}4$~m :
\[
  b = \frac{S_{\text{gaine}}}{a} = \frac{0{,}0444}{0{,}4} = 0{,}111~\text{m} \approx \boxed{110~\text{mm}}
\]

La section est $400 \times 110$~mm. La vitesse de 5~m/s est conforme aux recommandations pour une gaine principale tertiaire (limite : 6--8~m/s).

\textit{Question 4 — Réglages opérationnels réducteurs de bruit :}
\begin{itemize}
  \item Réduire la vitesse de rotation du ventilateur par le variateur de vitesse (VFD) : une réduction de 20\,\% de la vitesse diminue le bruit d'environ 6~dB.
  \item Diminuer le débit de soufflage aux heures creuses (occupation réduite) en ajustant les registres ou la consigne du variateur.
\end{itemize}
